% ══════════════════════════════════════════
%  CHAPITRE 3 — MÉTHODOLOGIE SCRUM
% ══════════════════════════════════════════
\chapter{Méthodologie Scrum et Spécifications Globales}

\section{Présentation de la démarche Scrum}
Le projet a été mené selon la méthodologie agile \textbf{Scrum}, structurée en deux sprints de durée égale. Cette approche itérative nous a permis de livrer un incrément fonctionnel à chaque fin de sprint et d'adapter les priorités en fonction des retours.

\begin{figure}[H]
\centering
\imgplaceholder{Figure~\ref{fig:scrum_cycle} — Cycle Scrum}{Diagramme illustrant le cycle Scrum avec Product Backlog, Sprint Planning, Daily Scrum, Sprint Review et Retrospective}
\caption{Cycle de développement Scrum appliqué au projet}
\label{fig:scrum_cycle}
\end{figure}

\textbf{Rôles~:}
\begin{itemize}
  \item \textbf{Product Owner}~: Mme Nesrine [Nom] (encadrante)
  \item \textbf{Scrum Master}~: Mohamed Chaari
  \item \textbf{Équipe de développement}~: Mohamed, Yessine, Moataz, Brahim
\end{itemize}

\section{Identification des acteurs}

\begin{table}[H]
\centering
\caption{Acteurs du système}
\label{tab:acteurs}
\begin{tabularx}{\textwidth}{lX}
\toprule
\textbf{Acteur} & \textbf{Description} \\
\midrule
Utilisateur    & Consulte le dashboard, les prévisions et les analyses \\
Administrateur & Gère l'infrastructure Cloud et les configurations \\
Analyste       & Consulte les analyses statistiques avancées et les alertes \\
Système        & Exécute automatiquement les pipelines et les DAGs \\
\bottomrule
\end{tabularx}
\end{table}

\section{Besoins fonctionnels et non fonctionnels globaux}

\subsection{Besoins fonctionnels}
\begin{enumerate}
  \item Collecter les données météo historiques (2020--2024) et temps réel pour 221 villes.
  \item Ingérer les données via un système de streaming fiable (Kafka).
  \item Stocker les données dans une base PostgreSQL Cloud (Supabase).
  \item Orchestrer les analyses automatisées (moyennes, corrélations, pics, tendances).
  \item Exposer les données via une API REST documentée.
  \item Afficher les résultats dans un dashboard interactif avec cartes et graphiques.
\end{enumerate}

\subsection{Besoins non fonctionnels}
\begin{itemize}
  \item \textbf{Performance}~: Ingestion de $>$200 enregistrements/batch sans perte.
  \item \textbf{Disponibilité}~: Infrastructure Cloud avec haute disponibilité (Supabase, Aiven).
  \item \textbf{Sécurité}~: Connexions SSL pour Kafka, mots de passe chiffrés.
  \item \textbf{Maintenabilité}~: Code modulaire avec tests unitaires ($>$100 tests).
  \item \textbf{Scalabilité}~: Architecture découplée permettant l'ajout de nouvelles villes ou sources.
\end{itemize}

\section{Diagramme de cas d'utilisation global}

\begin{figure}[H]
\centering
\imgplaceholder{Figure~\ref{fig:uc_global} — Diagramme de cas d'utilisation global}{Diagramme UML montrant les 4 acteurs (Utilisateur, Administrateur, Analyste, Système) et les cas d'utilisation principaux~: Consulter Dashboard, Voir Prévisions, Analyser Tendances, Recevoir Alertes, Configurer Infrastructure, Exécuter Pipeline}
\caption{Diagramme de cas d'utilisation global}
\label{fig:uc_global}
\end{figure}

\section{Product Backlog complet}

\begin{table}[H]
\centering
\caption{Product Backlog — Liste des User Stories}
\label{tab:backlog}
\begin{tabularx}{\textwidth}{clXcc}
\toprule
\textbf{US} & \textbf{Sprint} & \textbf{Description} & \textbf{Priorité} & \textbf{Points} \\
\midrule
US1 & S1 & Configuration de Supabase et Aiven            & Haute & 8 \\
US2 & S1 & Producer live et Consumer vers la DB            & Haute & 13 \\
US3 & S1 & Endpoints API /summary et /forecast             & Moyenne & 8 \\
US4 & S1 & Authentification JWT                             & Moyenne & 5 \\
US5 & S1 & Dashboard et graphiques de prévisions            & Haute & 13 \\
US6 & S1 & Validation des flux et plan de test              & Basse & 3 \\
\midrule
US7 & S2 & DAGs Airflow (moyennes, corrélations)            & Haute & 13 \\
US8 & S2 & Détection de pics (canicule, tempêtes)            & Haute & 8 \\
US9 & S2 & Documentation, monitoring et livraison finale     & Moyenne & 5 \\
\bottomrule
\end{tabularx}
\end{table}

\section{Architecture globale du système}

\begin{figure}[H]
\centering
\imgplaceholder{Figure~\ref{fig:archi_globale} — Architecture globale du pipeline}{Diagramme d'architecture montrant~: APIs Météo $\rightarrow$ Producer Python $\rightarrow$ Aiven Kafka $\rightarrow$ Consumer Python $\rightarrow$ Supabase PostgreSQL $\rightarrow$ Airflow DAGs $\rightarrow$ FastAPI $\rightarrow$ Dashboard Frontend}
\caption{Architecture globale du système Météo Tunisie}
\label{fig:archi_globale}
\end{figure}

\section{Planification des Sprints}

\begin{table}[H]
\centering
\caption{Répartition des Sprints}
\label{tab:sprints}
\begin{tabularx}{\textwidth}{lXl}
\toprule
\textbf{Sprint} & \textbf{Objectif} & \textbf{User Stories} \\
\midrule
Sprint 1 & Infrastructure Cloud et Flux Temps Réel    & US1 à US6 \\
Sprint 2 & Analyses Avancées, Alertes et Déploiement   & US7 à US9 \\
\bottomrule
\end{tabularx}
\end{table}
